% GENERATED FILE. Edit canonical lessons, then run npm run generate:books.
% canonical-source-hash: fnv1a64:fde6a5ffd5c0ba1d
% canonical-lessons: MW-C279-gothiyo, MW-C279-chaupar, MW-C279-paso, MW-C279-hod, MW-C279-kusti

\chapter{Games and Contests}
\label{ch:mw-games-and-contests}
\hlchaptermodality{279}

\begin{chapteropening}
\textbf{By the end of this chapter:} \emph{I can name the games I play with friends.}

Complete the last lesson of chapter 279: I can name the games I play with friends.
\end{chapteropening}

\section[\texorpdfstring{\mw{गोठियो} (goṭhiyo)}{गोठियो (goṭhiyo)}]{\mw{गोठियो} (goṭhiyo) — a friend; a companion}

\label{lesson:MW-C279-gothiyo}



\begin{quote}
Before the new one: say the Marwadi for a doll, then the Marwadi for to hide.
\end{quote}

\subsection*{You'll want to know: \mw{गोठियो}}
\textbf{\mw{गोठियो}} — \emph{goṭhiyo} — \textquotedblleft{}a friend; a companion\textquotedblright{}.

\begin{quote}
\mw{म्हारो} \mw{गोठियो} \mw{अठै} \mw{है।}
\end{quote}

\begin{quote}
\emph{mhāro goṭhiyo aṭhai hai.}
\end{quote}

\begin{quote}
My friend is here.
\end{quote}

A masculine noun.

\begin{cousinweb}
\textbf{\mw{गोठियो}} (\emph{goṭhiyo}) is inherited from Sanskrit \emph{goṣṭhika}, \textquotedblleft{}belonging to a gathering\textquotedblright{}; CDIAL gives no Marwari form in this entry, but lists a Gujarati form glossed \textquotedblleft{}friend\textquotedblright{} (CDIAL 4338a).
\end{cousinweb}

\begin{grammarlens}[title={What you've built}]
One more word for this chapter.
\end{grammarlens}

\subsection*{Your turn}
\begin{itemize}
\raggedright
  \item \emph{Say it:} \emph{goṭhiyo}
  \item \emph{Say it:} \emph{goṭhiyo}, once more
  \item \emph{Say it:} \emph{mhāro goṭhiyo aṭhai hai.} — My friend is here.
  \item \emph{Recall:} say the Marwadi for a doll, then the Marwadi for to hide, then say \emph{goṭhiyo} again
\end{itemize}

\begin{tcolorbox}[breakable,colback=teal!4,colframe=teal!35!black,title={Before you move on}]
What does \mw{गोठियो} mean? (\textquotedblleft{}A friend; a companion\textquotedblright{}.) Say it once more.
\end{tcolorbox}
\section[\texorpdfstring{\mw{चौपड़} (chaupaṛ)}{चौपड़ (chaupaṛ)}]{\mw{चौपड़} (chaupaṛ) — chaupar, a board game}

\label{lesson:MW-C279-chaupar}



\begin{quote}
Before the new one: say the Marwadi for to hide, then the Marwadi for a friend.
\end{quote}

\subsection*{You'll want to know: \mw{चौपड़}}
\textbf{\mw{चौपड़}} — \emph{chaupaṛ} — \textquotedblleft{}chaupar, a board game\textquotedblright{}.

\begin{quote}
\mw{दादा} \mw{चौपड़} \mw{खेलै} \mw{है।}
\end{quote}

\begin{quote}
\emph{dādā chaupaṛ khelai hai.}
\end{quote}

\begin{quote}
Grandfather plays chaupar.
\end{quote}

A feminine noun.

\begin{cousinweb}
\textbf{\mw{चौपड़}} (\emph{chaupaṛ}) is inherited from Sanskrit \emph{catuṣpuṭa}, \textquotedblleft{}having four folds\textquotedblright{}; CDIAL lists Marwari \emph{copṛī} and Old Marwari \emph{caupaḍī} for the game (CDIAL 4648), played on a cloth with four arms.
\end{cousinweb}

\begin{grammarlens}[title={What you've built}]
One more word for this chapter.
\end{grammarlens}

\subsection*{Your turn}
\begin{itemize}
\raggedright
  \item \emph{Say it:} \emph{chaupaṛ}
  \item \emph{Say it:} \emph{chaupaṛ}, once more
  \item \emph{Say it:} \emph{dādā chaupaṛ khelai hai.} — Grandfather plays chaupar.
  \item \emph{Recall:} say the Marwadi for to hide, then the Marwadi for a friend, then say \emph{chaupaṛ} again
\end{itemize}

\begin{tcolorbox}[breakable,colback=teal!4,colframe=teal!35!black,title={Before you move on}]
What does \mw{चौपड़} mean? (\textquotedblleft{}Chaupar, a board game\textquotedblright{}.) Say it once more.
\end{tcolorbox}
\section[\texorpdfstring{\mw{पासो} (pāso)}{पासो (pāso)}]{\mw{पासो} (pāso) — a die for games}

\label{lesson:MW-C279-paso}



\begin{quote}
Before the new one: say the Marwadi for a friend, then the Marwadi for a board game.
\end{quote}

\subsection*{You'll want to know: \mw{पासो}}
\textbf{\mw{पासो}} — \emph{pāso} — \textquotedblleft{}a die for games\textquotedblright{}.

\begin{quote}
\mw{म्हारो} \mw{पासो} \mw{लाल} \mw{है।}
\end{quote}

\begin{quote}
\emph{mhāro pāso lāl hai.}
\end{quote}

\begin{quote}
My die is red.
\end{quote}

A masculine noun.

\begin{cousinweb}
\textbf{\mw{पासो}} (\emph{pāso}) is inherited from Sanskrit \emph{pāśa}, \textquotedblleft{}a die\textquotedblright{}; CDIAL gives no Marwari form in this entry, but lists Hindi \emph{pāsā}, \textquotedblleft{}die\textquotedblright{} (CDIAL 8132).
\end{cousinweb}

\begin{grammarlens}[title={What you've built}]
One more word for this chapter.
\end{grammarlens}

\subsection*{Your turn}
\begin{itemize}
\raggedright
  \item \emph{Say it:} \emph{pāso}
  \item \emph{Say it:} \emph{pāso}, once more
  \item \emph{Say it:} \emph{mhāro pāso lāl hai.} — My die is red.
  \item \emph{Recall:} say the Marwadi for a friend, then the Marwadi for a board game, then say \emph{pāso} again
\end{itemize}

\begin{tcolorbox}[breakable,colback=teal!4,colframe=teal!35!black,title={Before you move on}]
What does \mw{पासो} mean? (\textquotedblleft{}A die for games\textquotedblright{}.) Say it once more.
\end{tcolorbox}
\section[\texorpdfstring{\mw{होड} (hoḍ)}{होड (hoḍ)}]{\mw{होड} (hoḍ) — a contest}

\label{lesson:MW-C279-hod}



\begin{quote}
Before the new one: say the Marwadi for a board game, then the Marwadi for a die for games.
\end{quote}

\subsection*{You'll want to know: \mw{होड}}
\textbf{\mw{होड}} — \emph{hoḍ} — \textquotedblleft{}a contest\textquotedblright{}.

\begin{quote}
\mw{आज} \mw{गांव} \mw{मांय} \mw{होड} \mw{है।}
\end{quote}

\begin{quote}
\emph{āj gā̃v mā̃y hoḍ hai.}
\end{quote}

\begin{quote}
There's a contest in the village today.
\end{quote}

A feminine noun.

\begin{cousinweb}
\textbf{\mw{होड}} (\emph{hoḍ}) is inherited from a reconstructed \emph{*hōḍḍa}, \textquotedblleft{}a wager\textquotedblright{}, which CDIAL takes from Dravidian; it lists Old Marwari \emph{hoḍa}, \textquotedblleft{}wager, rivalry\textquotedblright{} (CDIAL 14175).
\end{cousinweb}

\begin{grammarlens}[title={What you've built}]
One more word for this chapter.
\end{grammarlens}

\subsection*{Your turn}
\begin{itemize}
\raggedright
  \item \emph{Say it:} \emph{hoḍ}
  \item \emph{Say it:} \emph{hoḍ}, once more
  \item \emph{Say it:} \emph{āj gā̃v mā̃y hoḍ hai.} — There's a contest in the village today.
  \item \emph{Recall:} say the Marwadi for a board game, then the Marwadi for a die for games, then say \emph{hoḍ} again
\end{itemize}

\begin{tcolorbox}[breakable,colback=teal!4,colframe=teal!35!black,title={Before you move on}]
What does \mw{होड} mean? (\textquotedblleft{}A contest\textquotedblright{}.) Say it once more.
\end{tcolorbox}
\section[\texorpdfstring{\mw{कुस्ती} (kustī)}{कुस्ती (kustī)}]{\mw{कुस्ती} (kustī) — wrestling}

\label{lesson:MW-C279-kusti}



\begin{quote}
Before the new one: say the Marwadi for a die for games, then the Marwadi for a contest.
\end{quote}

\subsection*{You'll want to know: \mw{कुस्ती}}
\textbf{\mw{कुस्ती}} — \emph{kustī} — \textquotedblleft{}wrestling\textquotedblright{}.

\begin{quote}
\mw{भाई} \mw{कुस्ती} \mw{सीखै} \mw{है।}
\end{quote}

\begin{quote}
\emph{bhāī kustī sīkhai hai.}
\end{quote}

\begin{quote}
My brother is learning wrestling.
\end{quote}

A feminine noun.

\begin{cousinweb}
\textbf{\mw{कुस्ती}} (\emph{kustī}) is usually traced to Persian \emph{kushtī}, \textquotedblleft{}wrestling\textquotedblright{}.
\end{cousinweb}

\begin{grammarlens}[title={What you've built}]
One more word for this chapter.
\end{grammarlens}

\subsection*{Your turn}
\begin{itemize}
\raggedright
  \item \emph{Say it:} \emph{kustī}
  \item \emph{Say it:} \emph{kustī}, once more
  \item \emph{Say it:} \emph{bhāī kustī sīkhai hai.} — My brother is learning wrestling.
  \item \emph{Recall:} say the Marwadi for a die for games, then the Marwadi for a contest, then say \emph{kustī} again
\end{itemize}

\begin{tcolorbox}[breakable,colback=teal!4,colframe=teal!35!black,title={Before you move on}]
What does \mw{कुस्ती} mean? (\textquotedblleft{}Wrestling\textquotedblright{}.) Say it once more.
\end{tcolorbox}
